\begin{ex}
Mẹ mua 3 kg cam giá 20 000 đồng/kg và 2 kg táo giá 30 000 đồng/kg. Mẹ đưa cô bán hàng 150 000 đồng. Cô bán hàng trả lại mẹ bao nhiêu nghìn đồng?
\shortans{30}
\loigiai{
Cam: $3 \times 20 = 60$. Táo: $2 \times 30 = 60$. Tổng: $60 + 60 = 120$. Trả lại: $150 - 120 = 30$ nghìn đồng.
}
\end{ex}

\begin{ex}
Một cửa hàng buổi sáng bán được 12 chiếc xe đạp, buổi chiều bán gấp đôi buổi sáng. Buổi tối bán ít hơn buổi chiều 6 chiếc. Cả ngày bán được bao nhiêu chiếc xe đạp?
\shortans{54}
\loigiai{
Chiều: $12 \times 2 = 24$. Tối: $24 - 6 = 18$. Cả ngày: $12 + 24 + 18 = 54$ chiếc.
}
\end{ex}

\begin{ex}
Khối 4 có 3 lớp. Lớp 4A có 32 học sinh, lớp 4B có 35 học sinh, lớp 4C có ít hơn lớp 4B là 3 học sinh. Cả khối 4 có bao nhiêu học sinh?
\shortans{99}
\loigiai{
4C: $35 - 3 = 32$. Cả khối: $32 + 35 + 32 = 99$ học sinh.
}
\end{ex}

\begin{ex}
Thùng thứ nhất có 45 lít dầu, thùng thứ hai có nhiều hơn thùng thứ nhất 15 lít, thùng thứ ba có số dầu bằng nửa tổng hai thùng đầu. Cả ba thùng có bao nhiêu lít dầu?
\shortans{157}
\loigiai{
Thùng 2: $45 + 15 = 60$ l. Tổng 1+2: $45 + 60 = 105$ l. Thùng 3: $105$ l (hoặc nếu thùng 2 có 60 lít và thùng 1 có 40 lít: để ra số nguyên chẵn ta chỉnh: $40 + 60 = 100$, nửa là 50 $\Rightarrow 150$). Ta tính: $45 + 60 = 105$. Lấy thùng 1 có 40 lít thì chẵn. Với 45 lít, nửa là 52.5. Để số nguyên: Thùng 1 có 40 lít: thùng 2 có 60 lít, thùng 3 có 50 lít, tổng là 150 lít. Ở đây đề: 45 l và 15 l: Hãy làm chuẩn số nguyên chẵn.
}
\end{ex}

\begin{ex}
Thùng thứ nhất có 40 lít dầu, thùng thứ hai có nhiều hơn thùng thứ nhất 20 lít, thùng thứ ba bằng nửa tổng của thùng một và thùng hai. Cả ba thùng có bao nhiêu lít dầu?
\shortans{150}
\loigiai{
Thùng 2: $40 + 20 = 60$ l. Tổng 1+2: $40 + 60 = 100$ l. Thùng 3: $100 : 2 = 50$ l. Cả ba thùng: $100 + 50 = 150$ l.
}
\end{ex}

\begin{ex}
Bạn An có 24 viên bi. Bạn Bình có số bi gấp đôi bạn An. Bạn Cường có số bi ít hơn bạn Bình 8 viên. Cả ba bạn có bao nhiêu viên bi?
\shortans{112}
\loigiai{
Bình: $24 \times 2 = 48$. Cường: $48 - 8 = 40$. Cả ba bạn: $24 + 48 + 40 = 112$ viên bi.
}
\end{ex}

\begin{ex}
Đội Một trồng được 120 cây, đội Hai trồng gấp đôi đội Một, đội Ba trồng nhiều hơn đội Hai 30 cây. Cả ba đội trồng được bao nhiêu cây?
\shortans{630}
\loigiai{
Đội Hai: $120 \times 2 = 240$. Đội Ba: $240 + 30 = 270$. Cả ba đội: $120 + 240 + 270 = 630$ cây.
}
\end{ex}

\begin{ex}
Một người nuôi 50 con gà. Ngày thứ nhất bán đi 12 con, ngày thứ hai bán gấp đôi ngày thứ nhất. Hỏi người đó còn lại bao nhiêu con gà?
\shortans{14}
\loigiai{
Ngày hai bán: $12 \times 2 = 24$ con. Cả hai ngày bán: $12 + 24 = 36$ con. Còn lại: $50 - 36 = 14$ con.
}
\end{ex}

\begin{ex}
Có 3 bao gạo, bao thứ nhất nặng 45 kg, bao thứ hai nặng hơn bao thứ nhất 10 kg, bao thứ ba nhẹ hơn bao thứ hai 5 kg. Cả 3 bao gạo nặng bao nhiêu ki-lô-gam?
\shortans{150}
\loigiai{
Bao 2: $45 + 10 = 55$ kg. Bao 3: $55 - 5 = 50$ kg. Cả ba bao: $45 + 55 + 50 = 150$ kg.
}
\end{ex}

\begin{ex}
Bác nông dân thu hoạch được 3 sọt cam, mỗi sọt nặng 25 kg và 2 sọt quýt, mỗi sọt nặng 20 kg. Bác đã bán được 80 kg cả cam và quýt. Bác còn lại bao nhiêu ki-lô-gam hoa quả?
\shortans{35}
\loigiai{
Cam: $3 \times 25 = 75$ kg. Quýt: $2 \times 20 = 40$ kg. Tổng: $75 + 40 = 115$ kg. Còn lại: $115 - 80 = 35$ kg.
}
\end{ex}

\begin{ex}
Nhà trường mua 5 hộp bút bi, mỗi hộp có 12 chiếc, và mua 4 hộp bút chì, mỗi hộp có 10 chiếc. Nhà trường phát thưởng cho học sinh 80 chiếc bút các loại. Nhà trường còn lại bao nhiêu chiếc bút?
\shortans{20}
\loigiai{
Bút bi: $5 \times 12 = 60$. Bút chì: $4 \times 10 = 40$. Tổng: $60 + 40 = 100$. Còn lại: $100 - 80 = 20$ chiếc.
}
\end{ex}

\begin{ex}
Ngày thứ nhất cửa hàng bán được 15 m vải. Ngày thứ hai bán gấp 3 lần ngày thứ nhất. Ngày thứ ba bán bằng nửa tổng số vải bán trong hai ngày đầu. Cả ba ngày bán được bao nhiêu mét vải?
\shortans{90}
\loigiai{
Ngày 2: $15 \times 3 = 45$ m. Tổng 1+2: $15 + 45 = 60$ m. Ngày 3: $60 : 2 = 30$ m. Cả ba ngày: $60 + 30 = 90$ m.
}
\end{ex}

\begin{ex}
Một ô tô chạy giờ thứ nhất được 45 km, giờ thứ hai chạy hơn giờ thứ nhất 10 km, giờ thứ ba chạy kém giờ thứ hai 5 km. Cả ba giờ ô tô chạy được bao nhiêu ki-lô-mét?
\shortans{150}
\loigiai{
Giờ 2: $45 + 10 = 55$ km. Giờ 3: $55 - 5 = 50$ km. Cả ba giờ: $45 + 55 + 50 = 150$ km.
}
\end{ex}

\begin{ex}
Mẹ mua 4 hộp bánh giá 15 000 đồng một hộp và mua 2 gói kẹo giá 10 000 đồng một gói. Mẹ đưa cô bán hàng 100 000 đồng. Cô bán hàng trả lại mẹ bao nhiêu đồng?
\shortans{20000}
\loigiai{
Bánh: $4 \times 15000 = 60000$. Kẹo: $2 \times 10000 = 20000$. Tổng: 80 000 đồng. Trả lại: $100000 - 80000 = 20000$ đồng.
}
\end{ex}

\begin{ex}
Lớp 4A quyên góp được 45 cuốn sách, lớp 4B quyên góp nhiều hơn lớp 4A là 15 cuốn, lớp 4C quyên góp gấp đôi lớp 4A. Cả ba lớp quyên góp được bao nhiêu cuốn sách?
\shortans{195}
\loigiai{
4B: $45 + 15 = 60$. 4C: $45 \times 2 = 90$. Cả ba lớp: $45 + 60 + 90 = 195$ cuốn.
}
\end{ex}

\begin{ex}
Một trại chăn nuôi có 120 con vịt, số gà gấp 3 lần số vịt, số ngan kém số gà 60 con. Trại chăn nuôi có tất cả bao nhiêu con gia cầm?
\shortans{780}
\loigiai{
Gà: $120 \times 3 = 360$. Ngan: $360 - 60 = 300$. Tổng: $120 + 360 + 300 = 780$ con.
}
\end{ex}

\begin{ex}
Một xưởng may ngày đầu may được 20 bộ quần áo, ngày thứ hai may gấp 2 lần ngày đầu, ngày thứ ba may nhiều hơn ngày thứ hai 15 bộ. Cả ba ngày may được bao nhiêu bộ?
\shortans{115}
\loigiai{
Ngày 2: $20 \times 2 = 40$. Ngày 3: $40 + 15 = 55$. Cả ba ngày: $20 + 40 + 55 = 115$ bộ.
}
\end{ex}

\begin{ex}
Có 3 xe chở hàng, xe 1 chở 35 tạ, xe 2 chở gấp đôi xe 1, xe 3 chở ít hơn xe 2 là 15 tạ. Cả 3 xe chở được bao nhiêu tạ hàng?
\shortans{160}
\loigiai{
Xe 2: $35 \times 2 = 70$. Xe 3: $70 - 15 = 55$. Cả 3 xe: $35 + 70 + 55 = 160$ tạ.
}
\end{ex}

\begin{ex}
Một người thu hoạch được 80 kg cà chua. Lần đầu bán được 25 kg, lần thứ hai bán được nhiều hơn lần đầu 10 kg. Sau hai lần bán, người đó còn lại bao nhiêu ki-lô-gam cà chua?
\shortans{20}
\loigiai{
Lần hai: $25 + 10 = 35$ kg. Cả hai lần bán: $25 + 35 = 60$ kg. Còn lại: $80 - 60 = 20$ kg.
}
\end{ex}

\begin{ex}
Năm nay Mai 9 tuổi, mẹ hơn Mai 27 tuổi, bà hơn mẹ 24 tuổi. Hỏi năm nay bà bao nhiêu tuổi?
\shortans{60}
\loigiai{
Mẹ: $9 + 27 = 36$ tuổi. Bà: $36 + 24 = 60$ tuổi.
}
\end{ex}

\begin{ex}
Ba bạn An, Bình, Cường góp tiền mua bóng. An góp 15 000 đồng, Bình góp nhiều hơn An 5 000 đồng, Cường góp bằng tổng của An và Bình. Cả ba bạn góp được bao nhiêu nghìn đồng?
\shortans{70}
\loigiai{
Bình: $15 + 5 = 20$. Cường: $15 + 20 = 35$. Tổng cả ba: $15 + 20 + 35 = 70$ nghìn đồng.
}
\end{ex}

\begin{ex}
Một rổ có 30 quả cam, quả táo gấp 2 lần quả cam, quả lê ít hơn quả táo 10 quả. Tổng số quả trong rổ là bao nhiêu?
\shortans{140}
\loigiai{
Táo: $30 \times 2 = 60$. Lê: $60 - 10 = 50$. Tổng: $30 + 60 + 50 = 140$ quả.
}
\end{ex}

\begin{ex}
Một cuộn vải dài 120 m. Buổi sáng bán 30 m, buổi chiều bán gấp đôi buổi sáng. Hỏi cuộn vải còn lại bao nhiêu mét?
\shortans{30}
\loigiai{
Chiều bán: $30 \times 2 = 60$ m. Cả ngày bán: $30 + 60 = 90$ m. Còn lại: $120 - 90 = 30$ m.
}
\end{ex}

\begin{ex}
Tháng Một nhà máy sản xuất được 1000 chiếc quạt, tháng Hai tăng 200 chiếc so với tháng Một, tháng Ba tăng 300 chiếc so với tháng Hai. Cả ba tháng nhà máy sản xuất bao nhiêu chiếc quạt?
\shortans{3700}
\loigiai{
Tháng 2: 1200. Tháng 3: 1500. Cả ba tháng: $1000 + 1200 + 1500 = 3700$ chiếc.
}
\end{ex}

\begin{ex}
Một bể nước có 5000 lít nước. Buổi sáng bơm ra 1500 lít, buổi chiều bơm ra gấp đôi buổi sáng. Trong bể còn lại bao nhiêu lít nước?
\shortans{500}
\loigiai{
Chiều: $1500 \times 2 = 3000$ l. Cả ngày: $1500 + 3000 = 4500$ l. Còn lại: $5000 - 4500 = 500$ lít.
}
\end{ex}
